% Authoritative exercise source. Preserve all derivations and diagrams.

\ExerciseStart{G2.14}{Costruire uno stato composto e applicare una porta locale}{Collegamento}

\BookHeading{Problema, dati e obiettivo}

Preparare \(\ket a=(\ket0+i\ket1)/\sqrt2\) e \(\ket b=(\sqrt3\ket0+\ket1)/2\). Costruire \(\ket a\otimes\ket b\) nella base \((00,01,10,11)\) e applicare \(I\otimes H\). Calcolare tutte le ampiezze e probabilità finali.

\ExTheory{\ExRef{c2-tensore}{2.29}; \ExRef{c2-porte-locali}{2.30}}{Nel prodotto tensoriale \(c_{jk}=a_jb_k\). L’ordine della base è big-endian. La matrice \(I\otimes H\) applica H alla coppia \((00,01)\) e alla coppia \((10,11)\); lascia invariata l’etichetta del primo qubit.}

\ExPrediction{La porta agisce sul secondo fattore di una preparazione prodotto. Il primo qubit deve conservare le proprie statistiche, mentre le probabilità del secondo cambiano per interferenza.}

\ExFigure{\begin{tikzpicture}[>=Stealth,exlabel]
\draw[Ink] (0,1)--(5,1);\draw[Ink] (0,0)--(5,0);\node[left] at (0,1) {$\ket a$};\node[left] at (0,0) {$\ket b$};\node[draw=Accent,fill=white,minimum size=8mm] at (2.5,0) {$H$};\node[right] at (5,1) {primo qubit};\node[right] at (5,0) {secondo qubit};
\end{tikzpicture}}{La Hadamard è sul filo inferiore, che rappresenta il secondo fattore del prodotto tensoriale.}


\begin{ExTasks}

\item Verificare le norme dei due fattori.

\item Calcolare i quattro prodotti tensoriali.

\item Scrivere la matrice a quattro righe.

\item Eseguire tutti i prodotti riga-colonna.

\item Calcolare probabilità, marginali e fattorizzazione finale.

\end{ExTasks}

\ExSolution

\ExStep{Controllare i fattori}

\ExEq{\|a\|^2&=(1+|i|^2)/2=1,\\\|b\|^2&=(3+1)/4=1.}

\ExStep{Costruire le quattro ampiezze}

\ExEq{c_{00}&=\tfrac1{\sqrt2}\tfrac{\sqrt3}{2}=\tfrac{\sqrt3}{2\sqrt2},\quad
c_{01}=\tfrac1{\sqrt2}\tfrac12=\tfrac1{2\sqrt2},\\c_{10}&=\tfrac{i}{\sqrt2}\tfrac{\sqrt3}{2}=\tfrac{i\sqrt3}{2\sqrt2},\quad
c_{11}=\tfrac{i}{\sqrt2}\tfrac12=\tfrac{i}{2\sqrt2}.}
La norma totale è \((3+1+3+1)/8=1\).

\ExStep{Costruire la porta locale}

\ExEq{I\otimes H&=\frac1{\sqrt2}\begin{pmatrix}1&1&0&0\\1&-1&0&0\\0&0&1&1\\0&0&1&-1\end{pmatrix}.}
Le due copie di H agiscono nei sottospazi con primo bit fissato.

\ExStep{Moltiplicare ogni riga}

Il denominatore comune è \(\sqrt2\cdot2\sqrt2=4\):
\ExEq{c'_{00}&=(\sqrt3+1+0+0)/4,\\c'_{01}&=(\sqrt3-1+0+0)/4,\\c'_{10}&=(0+0+i\sqrt3+i)/4,\\c'_{11}&=(0+0+i\sqrt3-i)/4.}

\ExStep{Calcolare le probabilità}

\ExEq{P(00)=P(10)&=\frac{(\sqrt3+1)^2}{16}=\frac{2+\sqrt3}{8},\\P(01)=P(11)&=\frac{(\sqrt3-1)^2}{16}=\frac{2-\sqrt3}{8}.}
La somma è \(2[(2+\sqrt3)+(2-\sqrt3)]/8=1\). Le marginali del primo qubit sono \(1/2,1/2\); quelle del secondo sono \((2+\sqrt3)/4,(2-\sqrt3)/4\).

\ExStep{Verificare il prodotto finale}

\ExEq{H\ket b&=\frac1{2\sqrt2}\colv{\sqrt3+1\\\sqrt3-1},\\\ket a\otimes H\ket b&=\frac14\colv{\sqrt3+1\\\sqrt3-1\\i(\sqrt3+1)\\i(\sqrt3-1)}.}
La coincidenza con i quattro prodotti conferma quale filo sia stato trasformato.

\par\noindent\begin{minipage}{\linewidth}
\ExFigure{\begin{tikzpicture}[>=Stealth,exlabel]
\draw[exaxis] (-.6,0)--(4.4,0) node[right] {esito};\draw[exaxis] (-.6,0)--(-.6,3.6) node[above] {$p$};\node[left] at (-.6,3) {$ 1 $};\draw[Ink!15,dashed] (-.6,3)--(3.9000000000000004,3);\filldraw[fill=Accent!25,draw=Accent] (0.0,0) rectangle (0.65,1.399519052838329);\node[below] at (0.325,0) {$ 00 $};\node[above] at (0.325,1.399519052838329) {$ 0.4665 $};\filldraw[fill=Accent!25,draw=Accent] (1.1,0) rectangle (1.75,0.10048094716167105);\node[below] at (1.425,0) {$ 01 $};\node[above] at (1.425,0.10048094716167105) {$ 0.03349 $};\filldraw[fill=Accent!25,draw=Accent] (2.2,0) rectangle (2.85,1.399519052838329);\node[below] at (2.5250000000000004,0) {$ 10 $};\node[above] at (2.5250000000000004,1.399519052838329) {$ 0.4665 $};\filldraw[fill=Accent!25,draw=Accent] (3.3000000000000003,0) rectangle (3.95,0.10048094716167105);\node[below] at (3.6250000000000004,0) {$ 11 $};\node[above] at (3.6250000000000004,0.10048094716167105) {$ 0.03349 $};
\end{tikzpicture}}{Probabilità finali ordinate in base big-endian.}


\ExInterpret{La porta locale conserva la struttura prodotto. L’etichettatura dei fili, l’ordine delle componenti e la posizione dei blocchi nella matrice sono tre descrizioni della stessa operazione.}

\ExLab{Nelle sezioni 2.29–2.30 confrontare il prodotto tensoriale e la porta sul secondo sistema. Ripetere il calcolo con H sul primo filo usando \(H\otimes I\): cambiano le coppie di componenti mescolate. Il controllo della marginale identifica subito il qubit trasformato.}


\end{minipage}
